\documentclass[11pt]{article}
\usepackage[margin=1in]{geometry}
\usepackage[T1]{fontenc}
\usepackage{lmodern}
\usepackage{microtype}
\usepackage{booktabs}
\usepackage{graphicx}
\usepackage{amsmath}
\usepackage{xcolor}
\usepackage[hidelinks]{hyperref}
\usepackage{natbib}
\usepackage{enumitem}
\setlist{nosep,leftmargin=*}

\title{wev: Distilling LLM Browser Agents into Open, Local\\System-One Decision Models}
\author{AUTHOR NAME\\AFFILIATION\\\texttt{EMAIL}}
\date{}

\begin{document}
\maketitle

\begin{abstract}
System-One decision models answer typed questions (a choice, a yes/no, a score) about a free-form state with a
probability for every option, in one forward pass and without generating text. Browser agents make such decisions at
every step (\emph{which operation? which element?}), yet the open System-One models released so far are trained on
general classification and rule data. We show that they do not transfer to browser steps: on Mind2Web test requests
rendered in the request format of an open browser agent, Kev-4B and Kev-8B get 21\% and 19\% of steps fully right and
Laya under 1\%. We release \textbf{wev}, a family of open decision models (1.7B, 4B, 8B) built on Qwen3 base models,
trained on converted Mind2Web and NNetNav steps, general typed-decision corpora, and episodes collected from a
prompted LLM acting as the System One of a browser agent on live websites, verified by an LLM judge. wev-4b gets
75.9\% of Mind2Web test steps right while matching or exceeding Laya on general typed-decision benchmarks; it trails
Kev by 6--8 points on Kev's own suites. Used as the System One of the same agent on 153 held-out live-website tasks,
wev-4b completes 30 tasks and wev-8b 28, on par with the qwen3-max teacher (27), at 15--320\,ms per decision on one
consumer GPU. We also find that about a third of NNetNav's DONE labels do not hold up under an LLM judge. Code, data
builders, models and all evaluation outputs are public.
\end{abstract}

\section{Introduction}

Agents spend most of their steps on bounded decisions rather than on writing: which tool is next, whether a page shows
the answer, which element to click. System-One decision models \citep{jev} target exactly this. The caller declares
questions and their options at request time; the model returns a probability for every option, with nothing to parse
and nothing that can be malformed. Since the release of the hosted Jev model, open reproductions have appeared
\citep{kev,laya,thisthat}, together with studies of how to use such models inside agents \citep{reflex} and of their
failure modes \citep{typesafe}.

Browser agents are a natural client. The open agent jev-ultrafast \citep{jevultrafast} already asks a hosted System
One, at every step, which operation to perform and which listed element to target. Yet every open decision model we
know of is trained on general decision data. A browser step differs from those data in two ways: the state is long (a
page of text and 8--40 candidate elements, about 3{,}200 tokens in our data), and the right answer depends on the
page rather than on world knowledge or a rubric.

This report asks a practical question: \emph{can an open, local decision model serve as the System One of a browser
agent, and what does it take to train one?} Our contributions are:
\begin{itemize}
\item \textbf{Evidence that general decision models do not transfer to browser steps.} Scored on the same requests with
the same rule, Kev and Laya get at most 21\% of Mind2Web test steps right (\S\ref{sec:browser}).
\item \textbf{Data.} Converters from Mind2Web \citep{mind2web} and NNetNav \citep{nnetnav} to the agent's exact request
format; an LLM-judge relabelling of NNetNav's stopping decisions, which overturns 31\% of its DONE labels; and a
pipeline that collects judge-verified episodes from a prompted LLM running the agent on live sites (\S\ref{sec:data}).
\item \textbf{Models.} wev-1.7b, wev-4b and wev-8b: 76\% step success on Mind2Web test with general typed-decision
accuracy at or above Laya's, and live-website task success on par with the LLM teacher at 15--320\,ms per decision
(\S\ref{sec:results}).
\item \textbf{A reproducible evaluation.} One scorer for every model, per-request results, and an end-to-end harness
on live websites with judge-verified success. Everything is released.
\end{itemize}
We do not propose a new architecture: the model follows Kev's design (\S\ref{sec:model}), and two structural variants
we tried did not help at scale (\S\ref{sec:ablations}).

\section{Related work}

\paragraph{System-One decision models.} Jev \citep{jev} introduced the hosted model class; its training method (RLCD)
is unpublished. Kev \citep{kev} attaches LoRA adapters and a pointer head to Qwen base models and trains on ten public
classification sources plus rule records; Laya \citep{laya} fine-tunes ModernBERT-large \citep{modernbert} with a
small decision head; this-that-model \citep{thisthat} reads a 2B model's answer from a designated position. REFLEX
\citep{reflex} places hosted Jev inside an LLM agent and escalates to a strong LLM when confidence is low; Sun and Xu
\citep{typesafe} show that decision heads can follow the option name rather than the rubric bound to it. None of these
trains or evaluates a decision model on browser steps.

\paragraph{Web agents.} Mind2Web \citep{mind2web} provides human demonstrations on real websites; NNetNav
\citep{nnetnav} collects live-web trajectories by environment interaction. Most web agents generate actions with an
LLM. We instead cast each step as two typed questions and distil a generative teacher into a non-generative student,
in the spirit of knowledge distillation \citep{hinton2015} and imitation learning \citep{dagger}.

\section{Model}
\label{sec:model}

We follow Kev's design \citep{kev} and adapt its code (Apache-2.0).

\paragraph{Rendering.} A request is a state and a set of questions. The state is written once; each question follows
it as its own branch, with its options delimited by reserved tokens of the Qwen vocabulary:
\begin{center}\small
\texttt{<state> page text $\cdot$ elements $\cdot$ recent actions}\\
\texttt{<q> which operation? <opt> CLICK </opt> <opt> TYPE\_TEXT </opt> \ldots\ <decide>}\\
\texttt{<q> which element to CLICK? <opt> [1] Search </opt> <opt> [2] Sign in </opt> \ldots\ <decide>}
\end{center}
Delimiter strings that appear in user text are rewritten, so text cannot forge structure.

\paragraph{Mask and positions.} Every token of a question branch attends to the state and to its own branch, never to
another question; each branch restarts its positions after the state. One forward pass therefore answers all questions
of a request, and neither question order nor packing changes an answer (verified in tests to fp32 noise).

\paragraph{Backbone and readout.} A Qwen3 base model \citep{qwen3} with its vocabulary head removed (it cannot
generate), adapted with LoRA \citep{lora} (rank 16 on every attention and MLP projection). A pointer head projects the
hidden state of each option's \texttt{</opt>} token and of the question's \texttt{<decide>} token to 256 dimensions
and scores their dot product; a softmax over the options gives the answer. The adapter is merged into the released
weights.

\paragraph{Long states.} Browser states can exceed the context used in training. At inference the page text, URL,
per-element option lists and action history are shortened in that order until the request fits; a state of up to
4{,}096 tokens and question branches of up to 8{,}192 tokens are accepted.

\section{Data}
\label{sec:data}

Table~\ref{tab:data} lists the training mixture. All browser sources are rendered exactly as jev-ultrafast renders
its per-step request (state fields, operation options, instruction text, and \texttt{[index] label} targets).

\begin{table}[t]
\centering\small
\resizebox{\linewidth}{!}{%
\begin{tabular}{llrrl}
\toprule
source & license & train & weight & role \\
\midrule
Mind2Web \citep{mind2web} & CC BY 4.0 & 5{,}863 & 1 & human browser steps, split by website \\
NNetNav-live \citep{nnetnav}, relabelled & Apache-2.0 & 11{,}408 & 1 & live-web steps incl.\ DONE, BLOCKED, SCROLL \\
teacher episodes, first collection & qwen3-max outputs & 2{,}548 & 3 & judge-verified live-web episodes \\
teacher episodes, second collection$^\dagger$ & qwen3-max outputs & 2{,}274 & 3 & tasks the first collection failed \\
kev decision-v7 \citep{kev} & per source & 12{,}576 & 2 & classification, QA, rule records \\
typed-decisions (80\% of train) & Apache-2.0 & 960 & 8 & agent/ops workflows, 5 questions per case \\
tasksource-jev & mixed, some research-only & 57{,}024 & 1 & 392 classification tasks as decisions \\
jev-distill-corpus-v3 & Apache-2.0 & 50{,}000 & 1 & synthetic scenarios, soft labels \\
typed-decisions-synth & MIT & 6{,}682 & 1 & multi-question cases, 149 domains \\
\bottomrule
\end{tabular}}
\caption{Training mixture (rows after conversion; weight = repetitions per epoch). $^\dagger$wev-4b only.}
\label{tab:data}
\end{table}

\paragraph{Mind2Web.} Each recorded step becomes one request: the page's text and a sample of 8--40 candidate elements
(the gold element plus negatives), the task, and the preceding actions; the questions ask for the operation and, for
that operation, the target. Splits are by website, so test sites are unseen in training (train 5{,}863, dev 586, test
875 requests; 873 after removing two whose element lists exceed the API's 255-option limit).

\paragraph{NNetNav and its stopping labels.} NNetNav supplies the operations Mind2Web lacks: stopping (DONE), giving up
(BLOCKED) and scrolling. Its labels come from unsupervised exploration, and stopping labels are noisy. We asked an LLM
judge (DeepSeek-V4.1-Flash), given the goal and the page at each step, whether the goal was already achieved. On the
training split, the judge rejected 592 of the 1{,}898 DONE labels (31\%) and found the goal already achieved at 2{,}281
of the 10{,}102 other steps (23\%). We relabel the latter as DONE and drop the former, since the correct next step is
unknown.

\paragraph{Teacher episodes.} We serve a prompted LLM (qwen3-max) behind the System One API and let jev-ultrafast use
it on live websites, logging every request and answer. Tasks (1{,}402 train, 153 held-out evaluation) come from
NNetNav goals on their live sites, Wikipedia look-ups and Google Flights searches. We keep every decision of an episode
that ended in DONE when a judge, reading the final page, confirms the goal; every decision up to a BLOCKED on a page
that blocked the agent (CAPTCHA, 403, login wall); and, for episodes cut short by the browser infrastructure, every
decision but the last. Episodes that loop or exhaust their budget are dropped. Of 164 DONE claims in the second
collection, the judge confirmed 103. Train and dev are split by task.

We found two harness failures that inflated error rates in early collections: each episode left a browser daemon
running (thousands accumulated and starved the machine), and a page read during navigation aborted the episode. Both
are fixed in the released collector, and all end-to-end numbers below use the fixed harness.

\paragraph{General decisions.} Kev's decision-v7 training split, 80\% of the typed-decisions train split, and three
public corpora (tasksource-jev, jev-distill-corpus-v3, typed-decisions-synth), using soft labels where provided. Rows
whose state matches a typed-decisions test state are removed.

\section{Training}

One epoch, AdamW with learning rate $10^{-4}$ and a one-cycle schedule (10\% warm-up), cross-entropy against the
(soft) label distribution of every question, frozen bf16 base weights with fp32 adapters and head. Requests are packed
into length-bucketed micro-batches of at most 6{,}000 padded tokens per GPU for wev-4b and wev-8b, which train
data-parallel on 8 NVIDIA H20 GPUs (66 and 89 minutes), and of 20{,}000 tokens for wev-1.7b on one GPU (146 minutes).

\section{Evaluation}

\paragraph{Benchmarks.} Kev decision-v7 test (1{,}176 requests) is Kev's in-domain suite; Kev transfer-v4 test (764)
is out-of-domain for every model; typed-decisions test (400 cases, 5 questions each) covers four agent/ops workflows;
Mind2Web test (873) and NNetNav test (1{,}150) are browser steps.

\paragraph{Scoring.} Every model receives identical requests; its top option is its answer. We report per-question
accuracy for general benchmarks and \emph{step success} (operation and target both right) for browser steps. Kev runs
through its own server (bf16) and Laya through its Python package, both as published. Requests beyond wev's evaluation
context (11 of 873 on Mind2Web) count as wrong for wev.

\paragraph{End to end.} jev-ultrafast runs each of the 153 held-out tasks with the model as its System One; success
means the agent stops with DONE and a judge, reading the final page, agrees.

\paragraph{Test protocol.} Test splits were held out from training and model selection, with one exception: wev-4b
and wev-8b each had two final candidates (with and without the second teacher collection), and both were read on
test. The choice for wev-4b rests on development and end-to-end results; for wev-8b it was made after both test reads.

\section{Results}
\label{sec:results}

\subsection{General typed decisions}

\begin{table}[t]
\centering\small
\begin{tabular}{lccc}
\toprule
model & decision-v7 & transfer-v4 & typed-decisions \\
\midrule
wev-1.7b & 81.1 & 65.5 & \textbf{79.5} \\
wev-4b & 80.6 & 73.8 & 79.4 \\
wev-8b & 82.4 & 77.2 & 79.1 \\
\midrule
Kev-4B (Qwen3.5-4B) & \textbf{88.2} & \textbf{82.1} & 65.1 \\
Kev-8B (Qwen3-8B) & 88.1 & 76.8 & 62.7 \\
Laya (typed-decisions) & 65.7 & 62.8 & 76.8 \\
Laya & 64.3 & 63.7 & 36.2 \\
\bottomrule
\end{tabular}
\caption{General typed decisions, test splits, per-question accuracy (\%). wev and Laya (typed-decisions) train on the
typed-decisions train split; Kev and Laya do not.}
\label{tab:general}
\end{table}

Table~\ref{tab:general}. wev matches or exceeds both Laya models on every general benchmark. On Kev's own suites, Kev
leads: by 6--8 points on decision-v7, whose training split both train on, and on transfer-v4 Kev-4B leads by 8 points
over wev-4b while wev-8b is level with Kev-8B, which shares its base model. The typed-decisions column favours models
trained on that dataset: an earlier 4B wev trained without it (and without the external corpora) scored 53.1\% on its
development split. Our Laya (typed-decisions) score reproduces its published 76.6\%.

\subsection{Browser steps}
\label{sec:browser}

\begin{table}[t]
\centering\small
\begin{tabular}{lcc@{\hskip 2em}ccc}
\toprule
& \multicolumn{2}{c}{Mind2Web test} & \multicolumn{3}{c}{NNetNav test} \\
model & step & operation & step & DONE recall & premature DONE \\
\midrule
wev-1.7b & 68.2 & 88.1 & 61.0 & 80.4 & 9.4 \\
wev-4b & \textbf{75.9} & \textbf{91.2} & 61.4 & 84.6 & 10.0 \\
wev-8b & 75.5 & 90.3 & 60.8 & 79.6 & 8.2 \\
\midrule
Kev-4B & 21.2 & 35.7 & & & \\
Kev-8B & 19.0 & 73.3 & & & \\
Laya (typed-decisions) & 0.7 & 13.1 & & & \\
Laya & 0.0 & 2.5 & & & \\
\bottomrule
\end{tabular}
\caption{Browser steps (\%). Step: operation and target both right; targets are chosen among 8--40 candidates.
Premature DONE: share of not-finished steps answered DONE.}
\label{tab:browser}
\end{table}

Table~\ref{tab:browser}. General decision models fail on browser steps as published: Kev-8B picks the right operation
on 73\% of steps but the right element far less often, and Laya's context (512--1{,}024 tokens) holds only the start
of a page. wev gets three quarters of Mind2Web test steps right on websites it never saw. On NNetNav the hardest
decision is when to stop: wev recalls 80--85\% of DONE steps and says DONE too early on 8--10\% of the others;
thresholding the DONE probability trades one for the other.

\subsection{End to end on live websites}

\begin{table}[t]
\centering\small
\resizebox{\linewidth}{!}{%
\begin{tabular}{lcccc}
\toprule
System One & all (153) & NNetNav goals (123) & Wikipedia (10) & Google Flights (20) \\
\midrule
qwen3-max, prompted (teacher) & 27 & 13 & 10 & 4 \\
wev-4b & 30 & 18 & 10 & 2 \\
wev-8b & 28 & 18 & 10 & 0 \\
\bottomrule
\end{tabular}}
\caption{Judge-verified task successes on held-out live-website tasks with jev-ultrafast.}
\label{tab:e2e}
\end{table}

Table~\ref{tab:e2e}. The local students complete as many tasks as their teacher. The differences are within
run-to-run noise: two runs of wev-8b variants with equal totals disagreed on 8 tasks in each direction. Google Flights
often answers automated browsing with a CAPTCHA, and every model fails most of those tasks.

\subsection{Latency}

\begin{table}[t]
\centering\small
\begin{tabular}{lccc}
\toprule
& decision-v7 & typed-decisions & browser step \\
model & 1 q., $\sim$180 tok. & 5 q., $\sim$380 tok. & 2 q., $\sim$3{,}200 tok. \\
\midrule
wev-1.7b & 10\,ms & 13\,ms & 91\,ms \\
wev-4b & 15\,ms & 25\,ms & 217\,ms \\
wev-8b & 21\,ms & 37\,ms & 322\,ms \\
\bottomrule
\end{tabular}
\caption{Median in-process latency, one RTX 5090, bf16, one request at a time.}
\label{tab:latency}
\end{table}

Table~\ref{tab:latency}. All questions of a request share one forward pass, so five questions cost little more than
one; browser steps are dominated by encoding the page. On a 16\,GB Apple M5 laptop (PyTorch MPS backend, fp32),
wev-1.7b loads in about 30\,s and answers a one-question decision in 77\,ms.

\subsection{What did not help}
\label{sec:ablations}

\paragraph{Structural variants.} We tried two changes to the architecture on 1.7B (three seeds each, an earlier data
mix, development splits): a set head, which runs a two-layer permutation-equivariant transformer over the option
states before scoring, and truncating the backbone to its first 18 of 28 layers.

\begin{center}\small
\begin{tabular}{lcc}
\toprule
1.7B variant & Mind2Web dev step & NNetNav dev step \\
\midrule
pointer head, 28 layers & 69.0 $\pm$ 2.7 & 55.2 \\
set head, 28 layers & 66.5 & 53.9 \\
set head, 18 layers & 68.5 $\pm$ 0.2 & 56.6 \\
\bottomrule
\end{tabular}
\end{center}
The truncated model ran 23\% faster and was more stable across seeds, but a later same-hardware control run put its
cost at about 3 points, and the set head did not help at full depth. The released models use the pointer head and
every layer.

\paragraph{A second epoch.} Training wev-4b for two epochs changed development accuracy by less than one point on every
benchmark (transfer-v4 74.6$\to$75.1, decision-v7 81.9$\to$82.6, typed-decisions 79.7$\to$79.2, Mind2Web step
80.0$\to$79.3) at twice the cost.

\paragraph{The second teacher collection.} Adding episodes from tasks the first collection failed raised wev-4b's
live-website successes from 21 to 30; per task, 10 were gained and 1 lost. Its test accuracy was unchanged (within 0.6
points everywhere). On wev-8b the same data left live-website successes at 28 (8 gained, 8 lost) and lowered
transfer-v4 by 1.8 points, so wev-8b keeps the earlier mix.

\section{Limitations}

\begin{itemize}
\item \textbf{Evaluation size.} 153 live tasks give wide intervals, and live sites change between runs; we report
paired comparisons where we compare versions.
\item \textbf{Candidates.} Targets are chosen among the elements the agent lists (8--40 per step), not every element
on the page; the numbers are not comparable to the Mind2Web leaderboard.
\item \textbf{Test reads.} See the test protocol in \S5; wev-8b was chosen after two test reads.
\item \textbf{Scope.} English only; no comparison with hosted Jev (no API access); decisions only, with typed text
values produced by the agent's separate text model.
\item \textbf{Use terms.} Some training data is research-only (parts of tasksource-jev) or subject to its provider's
terms (qwen3-max outputs). We release the models as research artifacts.
\end{itemize}

\section{Conclusion}

General-purpose open decision models do not transfer to browser steps as published, but a decision model trained on
converted demonstrations, relabelled live-web data and judge-verified teacher episodes does: wev gets three quarters
of unseen-site steps right and, as the System One of an open browser agent, completes as many live tasks as the LLM it
was distilled from, locally and in tens to hundreds of milliseconds per decision. Natural next steps are on-policy
distillation, where the teacher labels the states the student reaches, and a larger live-task suite for tighter
end-to-end comparisons.

\paragraph{Release.} Code, data builders and evaluation outputs: \url{https://github.com/alanhuangyoo/wev}. Models:
\url{https://huggingface.co/alanhuangya}. Package: \texttt{pip install wev-ai}.

\bibliographystyle{plainnat}
\bibliography{refs}

\end{document}
